\chapter{Documentación del proyecto.}
\label{anexo:repos}

El código desarrollado durante el trabajo se reparte en tres repositorios de GitHub según la parte del trabajo a la que corresponde, de modo que cada uno de ellos puede consultarse.

\begin{itemize}
\item \textbf{Protocolo LoRaLAN}: \url{https://github.com/javicone/LoRaLAN_TFG}. Es el repositorio principal del trabajo, en el que se desarrolló el protocolo, y contiene la biblioteca LoRaLAN, el simulador de PC, el firmware de nodo, del LILYGO y monitor de re; el gateway que funciona sobre la Raspberry Pi y el servidor web. Cada directorio aporta su documentación.

\item \textbf{Sistema de control de asistencia}: 
\url{https://github.com/javicone/controlAsistenciaCUMe              }. Contiene el prototipo completo del sistema de fichaje en el que se basa el trabajo. Formado por el firmware del nodo de aula, el gateway para la Raspberry Pi, el servidor FastAPI y la base de datos, junto con las instrucciones de puesta en marcha.

\item \textbf{Medidas del espectro con HackRF}: \url{https://github.com/javicone/HackRF_sweep_setup}. Contiene las herramientas que fueron empleadas para operar el HackRF One y los scripts en Python que se utilizaron para analizar los resultados de la medición del espectro.
\end{itemize}
